The linear form $\tau_j^{\vert D\vert}$ does not vanish on commutators as could
be expected of a trace, leading to another discrepancy. The following result yields back  Proposition II.1 in~\cite{CoMo}.

\begin{proposition} Given two abstract pseudodifferential operators $A$, $B$ we have:
\begin{gather}\label{eq:taujDbracket}\tau_j^{D}([A,B])=\sum_{n= 1}^{k-j}
 (-1)^{n+1}\frac{\tau_{j+n}^{\vert D\vert} \left(
    A L^n(B)\right)}{n!},\qquad \forall\,
  j\in[[-1,k-1]],  \end{gather}
where we have set  $L(A):=  [\log \vert D\vert, A].$

Moreover $\tau_k([A,B])=0$ and when $k=0$ the above formula reads:
\[
\tau_{-1}^{D}([A,B ])=\tau_{0}\left( A \left([ \log
    \vert D\vert ,B] \right)\right).
    \]
\end{proposition}

\begin{proof}  Using the cyclic property of the
usual trace Tr, by analytic continuation we have:
\begin{gather*}
{\rm Tr} (B  A
  \vert D\vert^{-z} )\vert^{\rm mer} = {\rm Tr}\big(\big(\vert D\vert^{-z/2}B\big)  \big(A
  \vert D\vert^{-z/2}\big)\big)\vert^{\rm mer}\\
  \phantom{{\rm Tr} (B  A   \vert D\vert^{-z} )\vert^{\rm mer} }{}
 = {\rm Tr}\big( \big(A
  \vert D\vert^{-z/2}\big) \big( \vert D\vert^{-z/2}B\big)\big)\vert^{\rm mer}
 = {\rm Tr} ( A  \vert D\vert^{-z} B )\vert^{\rm mer}.
\end{gather*} Hence,
\begin{gather*}
{\rm Tr} ( A  B\vert D\vert^{-z}-B  A
  \vert D\vert^{-z} )\vert^{\rm mer} =  {\rm Tr} ( A  B \vert
  D\vert^{-z}-A \vert D\vert^{-z}B )\vert^{\rm mer}
 =  {\rm Tr} ( A [B ,\vert
  D\vert^{-z} ] )\vert^{\rm mer} .
\end{gather*}
We can apply Theorem \ref{thm:ResAz} to the holomorphic family
\[
A(z)= A\, [B, \vert
  D\vert^{-z}].
  \] Indeed,
\[
A(z) \vert
    D\vert^z= A  B -A \vert D\vert^{-z}B \vert
    D\vert^z= A   (B- \sigma(z)(B)),
    \]
where we have set $\sigma(z)(A)= \vert D\vert^{-z}  A \vert
    D\vert^z$. Since
    \[
    \vert D\vert^{-z} A\simeq  A  \vert D\vert^{-z}+
     \sum_{k=1}^\infty c_{-z,k}
\delta^k(A)   \vert D\vert^{-z-k},
\] it follows that
\[
\sigma(z)(A)\simeq A+ \sum_{k=1}^\infty c_{-z,k}
\delta^k(A)   \vert D\vert^{ -k},
\] where as before
 $c_{ \alpha, k}= \frac{\alpha(\alpha-1)\cdots (\alpha-k+1)}{k!}$ for any positive
 integer $k$. By
 Proposition \ref{prop:Aprime}, the  higher derivatives at zero $\partial_z^n (A(z) \vert
    D\vert^z)_{\vert_{z=0}}$  have order $a+b$, where $a$ is the order of $A$
    and $b$ that of $B$.

Since  $\partial^n_z \sigma(z)(A)_{\vert_{z=0}}=
(-1)^n\, L(A)$,   Theorem \ref{thm:ResAz}
    yields
\begin{gather*}
 \tau_j^{\vert D\vert}([A,B]) =  \sum_{n= 0}^{k-j} \frac{\tau_{j+n}^{\vert D\vert} \left(\partial_z^n (A(z)
  \vert D\vert^z)\right)_{\vert_{z=0}}}{n!}
 =  -\sum_{n= 0}^{k-j} \frac{\tau_{j+n}^{\vert D\vert} \left(A\, \left(\partial_z^n \sigma(z)
  (B)\right) \right)_{\vert_{z=0}}}{n!}\\
 \phantom{\tau_j^{\vert D\vert}([A,B])}{}
 =  \sum_{n= 0}^{k-j} (-1)^{n+1}\frac{\tau_{j+n}^{\vert D\vert} \left(
    A L^n(B)\right)}{n!}
 =  \sum_{n=1}^{k-j} (-1)^{n+1}\frac{\tau_{j+n}^{\vert D\vert} \left(
    A L^n(B)\right)}{n!}\qquad{\rm if}\quad j<k
    \end{gather*}
since the $n=0$ term vanishes.

  This computation also shows that $ \tau_k([A,B]) =\tau_k^{\vert D\vert} ([A,B])=0$.  When
  $k=0$ it  shows that:
\begin{gather*}
\tau_{-1}^{\vert D\vert}([A,B ])=\tau_{0}\left( A \left([ \log
    \vert D\vert ,B] \right)\right).\tag*{\qed}
    \end{gather*}\renewcommand{\qed}{}
\end{proof}

These discrepancy formulae are similar to the anomaly formulae for weighted
traces of classical  pseudodif\/ferential operators, see \cite{MN,CDMP,PS}  and
 \cite{Pa1}  for an extension  to  log-polyhomogeneous operators.


\section{ An analog of Kontsevich and Vishik's canonical trace}\label{section7}

In \cite{KV} Kontsevich and Vishik popularised what is known as the canonical trace on
non-integer order classical pseudodif\/ferential operators acting on smooth
sections of a vector bunlde $E$ over a closed manifold $
M$. This linear form
which vanishes on commutators of non-integer order classical pseudodif\/ferential
operators, is the unique (up to a multiplicative constant) linear extension of the ordinary trace to the set of
non-integer order operators with that property. The non integrality assumption
on the order can actually be weakened to  the order not lying in the set  $[{-}n,+\infty[\cap
\Z$ which corresponds to $- {\rm Sd}$, with ${\rm Sd}$ the dimension spectrum of the spectral triple
$\left(\Ci(M), L^2(M, E), D\right).$

For a spectral triple which fulf\/ills assumptions (H), (Dk) and (T), we replace
this  assumption on the order  with the requirement that the
order lies outside the set $-{\rm Sd}$.

\begin{lemma} For an abstract pseudodifferential  operator $A$ whose order~$a$ lies outside the set $ -{\rm Sd}$  and for any
  holomorphic family $A(z)$ such that $A(0)=A$,
 the map $z\mapsto {\rm Tr} \left( A(z) \right)\vert^{\rm mer}$ is holomorphic at
  $z=0$.

 In particular,  the map $z\mapsto  {\rm
    Tr} \left( A\, \vert D\vert^{-z}\right)\vert^{\rm mer} $ is holomorphic at
  zero so that
  \[
  \tau_j^{\vert D\vert}(A)=0\qquad \forall\, j\in \Z_+\qquad {\rm and}\qquad \tau_{-1}^{\vert D\vert} (P)=\lim_{z\to 0} ({\rm
    Tr}  ( A \vert D\vert^{-z} )\vert^{\rm mer} ).
    \]
\end{lemma}
\begin{proof} In view of the asymptotic expansion $A(z)\simeq \sum\limits_{j=0}^\infty
b_j (z)\vert D\vert^{\alpha(z)-j},$
where   $\alpha(z)$ is the  order of the holomorphic family $A(z)$,
it suf\/f\/ices to show the result for each holomorphic family $A_j(z)= b_j(z)
\vert D\vert^{\alpha(z)-j}$, with $j$ varying in $\Z_{\geq 0}$.
Such a family  has order $\alpha_j(z)= \alpha(z)-j$ so that the  poles of the meromorphic map $z\mapsto {\rm
    Tr} \left( A_j(z) \right)\vert^{\rm mer}$  lie in
  $\alpha^{-1}(-{\rm Sd}+\Z_{\geq 0})\subset \alpha^{-1}(-{\rm Sd})$, where we have used
  assumption (T) on the spectrum.
 Since    the order of $A$ does not lie in $-{\rm Sd}$, $0$ does not lie in this set of
poles, which leads to the two identities in the statement.\end{proof}
